\section{Introduccion}

En goofytex, la tarea se formula como retrieval visual multietiqueta: dada una imagen de bundle, el sistema debe recuperar y ordenar los \texttt{product\_asset\_id} mas probables, con un maximo de 15 predicciones por bundle.

El repositorio implementa un stack modular orientado a experimentacion rapida y despliegue reproducible:
\begin{enumerate}
  \item Preparacion de datos e imagenes (\texttt{download.sh}, \texttt{preprocess\_data.py}, \texttt{offline\_augment.py}).
  \item Entrenamiento de retrieval multimodal (\texttt{python -m src.train}, backend \texttt{src/models/retrieval\_openclip.py}).
  \item Inferencia y exportacion de submission (\texttt{src.infer}, \texttt{src.new\_infer}, \texttt{src.infer\_phase1}, \texttt{src.infer\_top1}, \texttt{src.infer\_openclip}).
\end{enumerate}

La documentacion adopta CRISP-DM como marco de trabajo para organizar decisiones tecnicas, trazabilidad y fases de iteracion. El objetivo de este paper es consolidar la metodologia y las tecnologias usadas en el codigo, sin incluir resultados de rendimiento.
